\documentclass[11pt]{article}

\usepackage[margin=1.1in]{geometry}
\usepackage{amsmath,amssymb,amsthm}
\usepackage{mathtools}
\usepackage{booktabs}
\usepackage{xcolor}
\usepackage{enumitem}
\usepackage[colorlinks=true,linkcolor=blue!60!black,citecolor=blue!60!black,urlcolor=blue!60!black]{hyperref}
\usepackage{microtype}

\newtheorem{theorem}{Theorem}[section]
\newtheorem{lemma}[theorem]{Lemma}
\newtheorem{definition}[theorem]{Definition}
\newtheorem{example}[theorem]{Example}
\newtheorem{proposition}[theorem]{Proposition}
\newtheorem{remark}[theorem]{Remark}

\newcommand{\vis}{\xrightarrow{\mathsf{vis}}}
\newcommand{\rc}{\xrightarrow{\mathsf{rc}}}
\newcommand{\comm}{\rightleftarrows}
\newcommand{\ncomm}{\not\rightleftarrows}
\newcommand{\Merge}{\mathsf{merge}}
\newcommand{\Sinit}{\sigma_0}
\newcommand{\dolab}{\mathsf{do}}
\newcommand{\Inv}{\mathsf{Inv}}
\newcommand{\app}{\mathsf{app}}
\newcommand{\commOn}[1]{\comm_{#1}}
\newcommand{\loE}[1]{\mathsf{lo}^{#1}}
\newcommand{\loA}[1]{\mathsf{lo}_{\!\app}^{#1}}
\newcommand{\eqobs}{\approx}
\newcommand{\resolve}{\mathsf{resolve}}
\newcommand{\anc}{\mathsf{anc}}
\newcommand{\live}{\mathsf{live}}
\newcommand{\reach}{\mathsf{Reach}}
\newcommand{\lean}[1]{\texttt{\small #1}}
\newcommand{\Ins}{\mathsf{Ins}}
\newcommand{\Del}{\mathsf{Del}}
\newcommand{\faith}{\mathsf{Faithful}}
\newcommand{\nfc}{\mathsf{NoFreshClash}}

\title{\bfseries A Verification Framework for State-Dependent MRDTs:\\
  \large describe only the sensible operations}
\author{Sal metatheory note --- framework draft\\
  \small FP Launchpad, IIT Madras}
\date{July 4, 2026}

\begin{document}
\maketitle

\begin{abstract}
\noindent
Mergeable replicated data types (MRDTs) are verified by a standard division of
labour: the data-type author discharges a fixed catalogue of local, equational
\emph{verification conditions} (VCs), and a once-and-for-all soundness
meta-theorem lifts them to RA-linearizability of every reachable configuration.
Every published catalogue quantifies its VCs \emph{over all states}. This is
adequate for \emph{flat} data types --- sets, counters, registers --- whose
operations are total: meaningful on every state. It is \emph{inadequate} for
\emph{state-dependent} data types, whose operations are only meaningful on
well-formed states and whose effect depends on the state they observe. The
sharp cases are tombstone-free sequence and rich-text CRDTs (a physical-delete
RGA; Peritext without tombstones): deletion physically removes a node and
\emph{re-anchors} its dependents, so an operation's recorded position can become
stale, and its commutation with a concurrent operation holds only on reachable
states. Whole-state VCs are simply \emph{false} on the infeasible states such a
type admits.

We present a \emph{conditioned} verification framework. The author enriches the
signature with two predicates --- an \emph{applicability} predicate $\app$
(``when is this operation sensible?'') and a reachability invariant $\Inv$ ---
and discharges VCs \emph{conditioned} on them; the framework proves
RA-linearizability once and for all. The developer thus \emph{describes only the
sensible operations} (as the RGA already does with insertion anchors) and is
never asked to reason about infeasible states. The technical core is an
\emph{applicability-aware} linearization order together with a commutation
discipline --- path-faithfulness plus causal freshness --- under which
state-dependent, re-anchoring operations provably commute. The development is
mechanized in Lean~4; we state which results are kernel-checked and which are
in progress.
\end{abstract}

\section{Introduction}
\label{sec:intro}

An MRDT is a replicated data type whose replicas fork, apply operations locally,
and reconcile divergence by a three-way \emph{merge} $\Merge(l,a,b)$ relative to
the lowest common ancestor $l$ of the two divergent states $a,b$. Correctness is
\emph{RA-linearizability}: every version in the store is explained by a
sequencing of exactly the events it has observed, consistent with a
linearization order $\mathsf{lo}$ derived from visibility and a
conflict-resolution relation $\mathsf{rc}$.

The prevailing methodology reduces this global property to a fixed catalogue of
\emph{local, equational} VCs over $\dolab$, $\Merge$ and $\mathsf{rc}$,
discharged per data type, and a soundness meta-theorem, proved once, that lifts
the VCs to RA-linearizability of every reachable configuration. Crucially, the
VCs are stated as universally quantified equations \emph{over all states}
$\sigma \in \Sigma$.

\paragraph{The whole-state assumption, and where it breaks.}
Universal quantification is benign when every operation is \emph{total} ---
meaningful on every state --- and \emph{state-independent} in the sense that a
VC true on the reachable states is true on all states. Flat data types are of
this kind: a grow-only set's $\mathsf{add}$ inserts on any state; an OR-set's
$\mathsf{rem}$ deletes whatever tags it sees. But consider a
\emph{tombstone-free} sequence CRDT (an RGA that physically deletes):

\begin{itemize}[leftmargin=1.4em]
\item An insertion records the \emph{position} it attaches to (a path to an
  anchor node). If a concurrent deletion has removed that anchor, the insertion
  must \emph{climb} to the nearest surviving ancestor --- its effect
  \emph{depends on the state}.
\item A deletion physically removes a node and \emph{re-anchors} (rehomes) its
  children to the deleted node's parent.
\item Consequently two operations commute only on \emph{well-formed} states:
  replayed on a syntactically-constructible but unreachable state, the equation
  is false.
\end{itemize}

A whole-state VC catalogue cannot express such a type: its commutation VC is
false as a $\forall\sigma$ statement, even though it holds on every state the
system can actually reach. This is not a peculiarity of the RGA; it is the
generic situation for any data type with \emph{structural} state and
\emph{positional} operations. Peritext without tombstones is the same
phenomenon one level up: a formatting mark anchors to character positions, and
deleting a character re-anchors the mark.

\paragraph{Thesis.}
Whole-state VC catalogues \emph{over-specify}: they demand equations on states no
execution reaches. The remedy is not to invent new equations but to
\emph{condition} the existing ones on what is reachable and sensible, and to
re-prove the meta-theorem to consume the conditioned forms. The author's added
burden should be exactly the \emph{description of sensible operations} --- which,
for the RGA, is already implicit in the anchor discipline --- and nothing more.

\paragraph{Contributions.}
\begin{enumerate}[leftmargin=1.6em]
\item A \emph{conditioned signature} (\S\ref{sec:sig}) carrying an applicability
  predicate $\app$ and a reachability invariant $\Inv$, with a \emph{collapse}
  property: flat types take $\app=\Inv=\top$ and recover the ordinary framework
  verbatim.
\item A \emph{conditioned VC discipline} (\S\ref{sec:vcs}): commutation and the
  linearization order are relativized to $\app$ and $\Inv$, so the author
  discharges VCs \emph{assuming sensibility}.
\item The technical core (\S\ref{sec:conv}): an \emph{applicability-aware}
  linearization order and a commutation discipline --- \emph{path-faithfulness}
  and \emph{causal freshness} --- under which state-dependent, re-anchoring
  operations commute; and the convergence theorem this yields.
\item A once-and-for-all \emph{adequacy} theorem (\S\ref{sec:adequacy}) reducing
  RA-linearizability to the conditioned VCs, with an explicit division between
  the author's obligations and the framework's generic proof.
\item Instantiation (\S\ref{sec:instances}): the tombstone-free RGA as the
  driving instance, and Peritext-without-tombstones as an instance of the same
  \emph{anchor-carrying} class that reuses the commutation discipline.
\end{enumerate}

\section{The replicated model and RA-linearizability}
\label{sec:model}

We recall the model, briefly, to fix notation; the framework of this note is
orthogonal to its details.

An \emph{event} is a triple $e=(t,r,o)$: a globally unique timestamp $t$, an
origin replica $r$, and an operation $o$. A \emph{configuration} carries a
\emph{store}: a directed acyclic graph of \emph{versions}, each holding a state
$\sigma\in\Sigma$ and the set $L(v)$ of events that produced it, with per-replica
heads. A transition forks a replica, applies a fresh event
($\sigma \mapsto \dolab(\sigma,e)$; we write $e(\sigma)$), queries, or
\textbf{merges}: it picks two heads $v_1,v_2$, finds their lowest common ancestor
$v_\ell$, and installs a new version with state
$\Merge(\sigma_{v_\ell},\sigma_{v_1},\sigma_{v_2})$ and event set
$L(v_1)\cup L(v_2)$. Visibility $\vis$ records what each event's issuing replica
had seen; $e_1 \parallel e_2$ denotes concurrency ($\neg(e_1\vis e_2)$ and
$\neg(e_2\vis e_1)$). Two events \emph{commute}, $e_1 \comm e_2$, when
$e_1(e_2(\sigma))=e_2(e_1(\sigma))$ for all $\sigma$.

\begin{definition}[RA-linearizability, per version]
\label{def:ralin}
A list $\pi=e^1\cdots e^n$ \emph{witnesses} a version holding state $s$ with event
set $E$ if (i) $\pi$ enumerates $E$ without repetition; (ii) $\pi$ extends the
linearization order $\mathsf{lo}$; and (iii) it folds to the state,
$e^n(\cdots e^1(\Sinit)\cdots)=s$. A configuration is \emph{RA-linearizable} if
every version in its store admits a witness.
\end{definition}

The order $\mathsf{lo}$ places $e_1$ before $e_2$ when either $e_1\vis e_2$ and
$e_1\ncomm e_2$, or $e_1\parallel e_2$, $e_1\ncomm e_2$, $\mathsf{rc}$ orders them,
and $e_2$ is not already overwritten by a later non-commuting event. We work
throughout with the \emph{set-relative} refinement $\loE{E}$ of this order, in
which the overwriter clause ranges over the event set $E$ of the version being
linearized rather than the whole configuration; this refinement is forced (an
event can gain overwriters from a concurrent branch) and $\mathsf{lo}\subseteq
\loE{E}$ pointwise, so every $\loE{E}$-witness is an $\mathsf{lo}$-witness.

\section{The conditioned signature}
\label{sec:sig}

\begin{definition}[conditioned MRDT signature]
\label{def:sig}
A \emph{conditioned MRDT signature} is a tuple
\[
  \mathcal{D} \;=\;
  \langle\, \Sigma,\ \Sinit,\ \dolab,\ \Merge,\ \mathsf{rc},\
  \eqobs,\ \Inv,\ \app \,\rangle
\]
extending an ordinary MRDT signature $\langle
\Sigma,\Sinit,\dolab,\Merge,\mathsf{rc}\rangle$ with
\begin{itemize}[leftmargin=1.4em]
\item an \emph{observational equivalence} $\eqobs\ \subseteq \Sigma\times\Sigma$
  --- an equivalence relation that is a congruence for $\dolab$ and $\Merge$
  (Def.~\ref{def:obs}); on flat types it may be taken as structural equality;
\item a \emph{reachability invariant} $\Inv:\Sigma\to\mathsf{Prop}$ --- an
  over-approximation of the reachable states;
\item an \emph{applicability} predicate $\app:\mathsf{Op}\times\Sigma\to
  \mathsf{Prop}$ --- ``operation $o$ is sensible on state $\sigma$.''
\end{itemize}
A signature is \emph{flat} when $\Inv=\top$, $\app=\top$, and $\eqobs$ is
structural equality; then all conditioned notions below reduce definitionally to
the ordinary ones.
\end{definition}

\begin{definition}[observational equivalence]
\label{def:obs}
$\sigma_1 \eqobs \sigma_2$ iff every query supported by the data type returns the
same value on $\sigma_1$ and $\sigma_2$. Convergence is stated up to $\eqobs$:
two replicas that have observed the same events must be observationally
equivalent, not necessarily structurally equal.\footnote{This is the notion of
\emph{convergence modulo observable behaviour}: two states may differ in
representation (a balanced vs.\ unbalanced search tree; a physically-deleted node
left as unreachable junk) yet agree on every observation. For the tombstone-free
RGA it turns out that reachable convergence holds even \emph{structurally}, so
$\eqobs$ may be taken as equality; we keep it abstract because state-dependent
types in general need it.}
\end{definition}

\paragraph{Running example: the tombstone-free RGA.}
The state is a forest over node identifiers: $\Sigma$ maps each live identifier
$t$ to a payload and an anchor $\anc(t)$ (its parent), with a root sentinel $0$.
We write $\live(\sigma,t)$ for ``$t$ is present in $\sigma$''. Two operations:
\begin{align*}
  \Ins(e,p,a)&:\quad \dolab(\sigma,\Ins(e,p,a)) = \sigma\big[t \mapsto (e,\
    \resolve(\sigma, a\!:\!p))\big], \\
  \Del(p,x)&:\quad \text{rehome each child of $x$ to } \resolve(\sigma,p),\
    \text{then remove } x,
\end{align*}
where $\resolve(\sigma,L)$ returns the first identifier of the list $L$ that is
live in $\sigma$, or the root $0$ if none is (the \emph{climb}). An insertion
records a \emph{path} $p$ and an \emph{anchor} $a$; a deletion records the path
$p$ to which its children rehome. The applicability and invariant predicates
capture ``sensible'':
\begin{align*}
  \app(\Ins(e,p,a),\sigma) &\;\equiv\; \big(a=0 \wedge p=[\,]\big)\ \vee\
    \big(\live(\sigma,a)\wedge \text{$p$ is the true ancestor chain of $a$}\big),\\
  \app(\Del(p,x),\sigma) &\;\equiv\; \text{analogously, $x$ live and $p$ its
    true rehome path},\\
  \Inv(\sigma) &\;\equiv\; \mathsf{wf}(\sigma)\ \wedge\ \mathsf{idmono}(\sigma),
\end{align*}
where $\mathsf{wf}$ says the forest is well-formed (every anchor live or the
root) and $\mathsf{idmono}$ says anchors strictly decrease
($\anc(\sigma,t)=0 \vee \anc(\sigma,t)<t$), a consequence of monotone timestamp
allocation. The anchor discipline the RGA already carries \emph{is} its
applicability predicate; the author writes nothing new to describe sensible
operations.

\begin{example}[the unconditioned commutation VC is false]
\label{ex:uncond}
This is why conditioning is necessary, not merely convenient. Take two
insertions issued concurrently by different replicas: $a$ records anchor $3$ and
path $[7]$ (it creates a node under $3$, climbing the path if $3$ is gone), and
$b$ creates a node with identifier $7$ at the root. Evaluate the two orders on the
initial root-only state $\Sinit$:
\begin{itemize}[leftmargin=1.4em]
\item $\dolab(\dolab(\Sinit,a),b)$ --- $a$ first: it resolves its list $[3,7]$
  against $\Sinit$, where both $3$ and $7$ are absent, and \emph{climbs to the
  root}, anchoring its node at $0$; then $b$ creates node $7$.
\item $\dolab(\dolab(\Sinit,b),a)$ --- $b$ first: it creates node $7$; now $a$
  resolves $[3,7]$ with $3$ absent but $7$ \emph{live}, anchoring its node at $7$.
\end{itemize}
The two results disagree on the parent of $a$'s node ($0$ versus $7$): the
whole-state commutation equation
\[
  \forall\sigma.\qquad \dolab(\dolab(\sigma,a),b)\ =\ \dolab(\dolab(\sigma,b),a)
\]
is therefore \emph{false}, already at $\sigma=\Sinit$
(\lean{naive\_general\_swap\_false}).
No execution reaches this counterexample: $a$ records an anchor $3$ that is not
present and a path naming $b$'s identifier $7$, yet $a$ and $b$ are
\emph{concurrent} --- $a$'s replica never observed $b$, so it could neither anchor
to an absent node nor name $b$'s fresh identifier. These are exactly the tuples
$\app$ (and causal freshness, \S\ref{sec:conv}) excludes; restricted to applicable
tuples, the two insertions \emph{do} commute (Lemma~\ref{lem:swap}). A whole-state
catalogue quantifies its commutation VC over $\sigma=\Sinit$ with this $a$ and so
cannot state a true VC for the type at all; the conditioned catalogue can. The RGA
already encodes the fix --- its shipped commutation VC is \emph{gated} on
\lean{accurate}${}\wedge{}$\lean{fresh}${}\wedge{}$\lean{wf} precisely because the
ungated equation is this counterexample.
\end{example}

\section{Sensible operations and reachability}
\label{sec:sensible}

The two predicates play complementary roles, and each carries exactly one
obligation.

\begin{definition}[reachability invariant]
\label{def:reach}
$\Inv$ is a \emph{reachability invariant} if $\Inv(\Sinit)$ holds and $\Inv$ is
preserved by every transition: for all $\sigma$ with $\Inv(\sigma)$ and every
applicable event $e$, $\Inv(e(\sigma))$; and for all $\Inv$-states $l,a,b$
arising as an honest LCA-and-branches triple, $\Inv(\Merge(l,a,b))$. We write
$\reach$ for the states satisfying $\Inv$.
\end{definition}

Establishing $\Inv$ as a reachability invariant is the author's obligation; it is
a routine induction over $\dolab$ and $\Merge$ (for the RGA, preservation of
$\mathsf{wf}$ under $\Merge$ requires exactly $\mathsf{idmono}$, which the same
induction supplies). No other whole-state reasoning is asked of the author.

\begin{remark}[the developer-effort ledger]
\label{rem:ledger}
The essential/author column and the generic/framework column are:
\begin{center}\small
\begin{tabular}{@{}p{0.44\textwidth}p{0.44\textwidth}@{}}
\toprule
Author provides (essential) & Framework proves (once) \\
\midrule
$\app$ --- spec of sensible operations & convergence of concurrent operations \\
$\Inv$ + reachability-invariance (Def.~\ref{def:reach}) & existence of canonical
  states $\sigma(E)$ \\
conditioned VCs (\S\ref{sec:vcs}, assuming $\app$) & the linearization machinery \\
& RA-linearizability (Thm.~\ref{thm:adequacy}) \\
\bottomrule
\end{tabular}
\end{center}
The author's column is the \emph{specification} --- one cannot verify a
state-dependent type without saying which states are valid and which operations
sensible --- and no more. The failure mode of a whole-state framework is that it
leaks framework-column work into the author's lap; \S\ref{sec:conv} shows how the
conditioned framework keeps the columns separate.
\end{remark}

\section{The conditioned verification conditions}
\label{sec:vcs}

The single primitive is conditioned commutation.

\begin{definition}[conditioned commutation]
\label{def:commOn}
Operations $o_1,o_2$ \emph{commute where sensible}, written $o_1 \commOn{\app}
o_2$, if for every $\sigma$ with $\Inv(\sigma)$, $\app(o_1,\sigma)$ and
$\app(o_2,\sigma)$,
\[
  \dolab(\dolab(\sigma,o_1),o_2)\ \eqobs\ \dolab(\dolab(\sigma,o_2),o_1).
\]
\end{definition}

Commutation is demanded only at reachable states where both operations are
sensible, and only up to observation. The relational VCs on $\mathsf{rc}$ are the
usual ones with $\comm$ replaced by $\commOn{\app}$: \textsc{rc-non-comm}
(non-commuting distinct-replica operations are $\mathsf{rc}$-ordered in exactly
one direction), \textsc{no-rc-chain} (no $o_1\rc o_2 \rc o_3$), and a conditioned
\textsc{cond-comm} (a resolved conflict is invisible once an overwriter runs).
These are the author's update-layer obligations, discharged \emph{assuming}
applicability of the operations involved --- which, for state-dependent types, is
what makes them provable at all.

\section{Canonical states and convergence}
\label{sec:conv}

This section is the technical core: how conditioned VCs yield a well-defined
canonical state per event set, and the commutation discipline that makes it work
for re-anchoring operations.

\subsection{The applicability-aware order}

The set-relative order $\loE{E}$ (\S\ref{sec:model}) orders \emph{non-commuting}
pairs. For state-dependent operations a second source of ordering appears: one
operation's \emph{applicability} may depend on another. A concurrent deletion can
render an insertion's recorded anchor stale; the insertion still commutes with
the deletion (it climbs), but the two are related by a \emph{generation
dependency} that the order must record, lest a reordering apply an operation at a
state where it is neither sensible nor a harmless no-op.

\begin{definition}[generation dependency; applicability-aware order]
\label{def:loA}
Event $e_2$ \emph{depends on} $e_1$, written $e_1 \prec_{\app} e_2$, if $e_1$ can
change whether $e_2$ is applicable: $\exists\sigma.\ \app(e_2,\sigma)\neq
\app(e_2,e_1(\sigma))$. The \emph{applicability-aware order} is
\[
  \loA{E} \;=\; \loE{E}\ \cup\ \{(e_1,e_2)\ :\ e_1\vis e_2\ \wedge\ e_1
    \prec_{\app} e_2\}.
\]
\end{definition}

For the RGA, $\prec_{\app}$ is the syntactic fact ``$e_2$'s recorded path names
$e_1$'s identifier''; it is discharged directly from the operation shapes.

\begin{definition}[no-op-feasible enumeration]
\label{def:noop}
A list $\pi=e^1\cdots e^n$ is \emph{no-op-feasible} if at each prefix fold
$\sigma_i = e^i(\cdots e^1(\Sinit)\cdots)$ the next event is either applicable at
$\sigma_i$ or acts as the identity there ($e^{i+1}(\sigma_i)\eqobs \sigma_i$).
\end{definition}

Two feasibility repairs are needed and are \emph{complementary}, not
alternatives: $\loA{E}$ excludes reorderings that violate a generation dependency
(which mere no-op tolerance would wrongly admit), while no-op tolerance admits the
harmless \emph{redundant} reorderings that strict applicability would wrongly
reject (e.g.\ two concurrent deletions of one node: whichever runs second is a
no-op, so no fully-applicable enumeration exists, yet both orders converge). The
order handles the first; no-op-feasibility the second.

\subsection{The commutation discipline for anchor-carrying types}

The convergence proof reorders an enumeration to a canonical form by adjacent
swaps of $\loA{E}$-incomparable events. Each swap must hold \emph{at the prefix
fold where the pair sits} --- a state at which either event may be stale, because
prior events in the prefix may have deleted its anchor. The discipline below
identifies exactly what makes such a swap valid without demanding either event be
applicable there.

\begin{definition}[path-faithfulness]
\label{def:faith}
An operation is \emph{faithful} at $\sigma$ if its recorded list tracks the true
tree of $\sigma$ above its climb target: for an insertion with list $L$, writing
$c=\resolve(\sigma,L)$, if $c$ is live then
$\resolve(\sigma,L\!\setminus\!\{c\}) = \anc(\sigma,c)$ (the list continues up the
real ancestor chain); a deletion is faithful if its rehome path likewise resolves
to its target's true parent. Write $\faith(o,\sigma)$.
\end{definition}

\begin{definition}[causal freshness]
\label{def:nfc}
$\nfc(a,b)$ holds if $b$'s (fresh) identifier does not occur in $a$'s recorded
list --- the condition a real execution guarantees between concurrent events,
since neither observed the other and identifiers are allocated afresh.
\end{definition}

Faithfulness is strictly weaker than applicability ($\app(o,\sigma)\Rightarrow
\faith(o,\sigma)$) and, unlike applicability, is \emph{stable under
re-anchoring}: a deletion that rehomes nodes preserves the faithfulness of a
concurrent operation's recorded path, because rehoming updates the true tree
consistently with the climb. This is the property that lets the swap fire at
stale states.

\begin{lemma}[faithful commutation]
\label{lem:swap}
Let $\sigma$ satisfy $\Inv$ with root absent, and let $a,b$ be distinct-timestamp
events with $\faith(a,\sigma)$, $\faith(b,\sigma)$, $\nfc(a,b)$ and $\nfc(b,a)$.
Then
\[
  \dolab(\dolab(\sigma,a),b)\ \eqobs\ \dolab(\dolab(\sigma,b),a).
\]
\hfill\lean{general\_swap\_bothFaithful}
\end{lemma}

Lemma~\ref{lem:swap} is mechanized (kernel-checked, all four
insertion/deletion combinations). It is the reusable heart of the framework for
anchor-carrying types: it says two re-anchoring operations commute whenever each
is merely faithful and they do not name each other's fresh identifiers ---
neither need be applicable. The asymmetric special case ($b$ applicable, $a$
faithful) is \lean{general\_swap}; that the symmetric case also holds is what
guarantees the convergence swap fires \emph{at every} prefix fold, where both
events may be stale.

\subsection{Convergence}

Faithfulness must hold not merely at one state but at every prefix fold the
canonicalization visits. This is a reachability-style invariant --- call it
\emph{faithful-threading} --- and it is where the framework absorbs the work that
a whole-state approach would push onto the author.

\begin{proposition}[faithful-threading, target form]
\label{prop:thread}
Along any $\loA{E}$-respecting, no-op-feasible enumeration of a backward-closed
$E$, every event is faithful at the prefix fold at which it is applied.
\end{proposition}

Threading is established from the author's $\app$/$\Inv$ data by a single-step
preservation argument: faithfulness is preserved by an applicable insertion, and
by an applicable deletion (the deletion splices its target out of the true chain,
re-linking children to the target's successor). The insertion step and the
correct invariant strengthening are mechanized; the deletion step is the one
remaining lemma (\S\ref{sec:status}).

\begin{theorem}[conditioned convergence, target form]
\label{thm:conv}
Assume the conditioned VCs (\S\ref{sec:vcs}), Lemma~\ref{lem:swap}, and
Proposition~\ref{prop:thread}. Then any two enumerations of a backward-closed
event set $E$ that respect $\loA{E}$ and are no-op-feasible fold from $\Sinit$
to $\eqobs$-equivalent states.
\end{theorem}

Convergence licenses the \emph{canonical state} $\sigma(E)$ --- the
$\eqobs$-class obtained by folding any such enumeration --- and reduces
``a version linearizes'' to ``its stored state $\eqobs\ \sigma$ of its event
set,'' with witness lists eliminated from all downstream reasoning. The proof is
the ordinary bottom-up canonicalization (proved once for flat types), with each
adjacent swap discharged by Lemma~\ref{lem:swap} under
Proposition~\ref{prop:thread}; nothing in it mentions $\resolve$, anchors, or the
RGA. That is the sense in which convergence is generic: it consumes the
faithful-commutation VC (Lemma~\ref{lem:swap}) and the threading obligation
(Proposition~\ref{prop:thread}) as hypotheses, and each data type supplies those
from its own operation algebra.

\section{Adequacy}
\label{sec:adequacy}

Convergence is the update-layer half. The merge-layer half is a single equation
between canonical states.

\begin{definition}[Join Lemma]
\label{def:join}
For backward-closed event sets $E_1,E_2$ of a reachable configuration,
$\sigma(E_1\cup E_2)\ \eqobs\ \Merge\big(\sigma(E_1\cap E_2),\ \sigma(E_1),\
\sigma(E_2)\big)$.
\end{definition}

By the LCA property $L(v_\ell)=L(v_1)\cap L(v_2)$, the first argument is exactly
the state the store hands the merge. The two halves compose once and for all:

\begin{theorem}[adequacy]
\label{thm:adequacy}
If $\mathcal{D}$ satisfies the conditioned VCs, has $\Inv$ as a reachability
invariant, and satisfies the Join Lemma, then every configuration reachable in
the transition system is RA-linearizable, per version.
\end{theorem}

The adequacy bridge is closure-indexed and proved generically: because every
store version is fully causally closed, the Join Lemma applies at it regardless
of the closure strength a data type declares, and the two established
verification routes (a delta contract for set- and counter-shaped merges; a
direct join for counter-comparison merges) are its two instances. The bridge and
its instances are mechanized (\lean{ra\_linearizable3\_of\_joinC}), with the nine
flat production MRDTs of the repository routed through it.

\section{The class map and instances}
\label{sec:instances}

Data types occupy a three-axis space: the \emph{delta class} of the merge (group,
lattice, or feasibility-conditioned), the \emph{closure strength} the Join Lemma
needs, and the \emph{state invariant} $\Inv$. Flat types sit on the $\Inv=\top$
plane and are handled by the ordinary framework (the collapse of
Def.~\ref{def:sig}). The conditioned framework of this note is what populates the
$\Inv\neq\top$ plane.

\paragraph{Tombstone-free RGA.} The driving instance: $\app$ = the anchor
discipline, $\Inv=\mathsf{wf}\wedge\mathsf{idmono}$, faithful commutation
(Lemma~\ref{lem:swap}) discharged from $\resolve$/rehome algebra.

\paragraph{Peritext without tombstones.} The same class, one level up. A
formatting mark anchors to character positions as an RGA node anchors to its
parent; deleting a character re-anchors the marks that span it, exactly the
rehoming phenomenon. The author supplies $\app$ (a mark's endpoints anchor live
characters), $\Inv$ (well-formed marks over a well-formed sequence), and
discharges the conditioned VCs by the \emph{same} faithfulness/freshness
discipline (Lemmas~\ref{lem:swap} and the threading of
Prop.~\ref{prop:thread}), with marks reusing the anchor machinery. It is thereby
an \emph{instance}, not a second bespoke development --- which is the payoff of a
framework over a per-type proof.

\section{Mechanization status}
\label{sec:status}

The development is in Lean~4 (Mathlib; axioms $\{\mathtt{propext},
\mathtt{Classical.choice}, \mathtt{Quot.sound}\}$; no \texttt{sorry} in the cited
results). Honestly delineated:

\begin{itemize}[leftmargin=1.4em]
\item \textbf{Kernel-checked.} The closure-indexed adequacy bridge
  (\lean{ra\_linearizable3\_of\_joinC}) and the nine flat MRDTs routed through it;
  the set-relative convergence for flat types; the conditioned-convergence order
  reduction; and the commutation discipline for the RGA ---
  Lemma~\ref{lem:swap} (\lean{general\_swap\_bothFaithful}) and its asymmetric
  form (\lean{general\_swap}), all insertion/deletion combinations.
\item \textbf{In progress.} Proposition~\ref{prop:thread}
  (faithful-threading): the correct invariant (a chain-faithfulness
  strengthening) and its insertion step are mechanized; the deletion-preservation
  step is the single remaining lemma, being discharged concurrently by direct
  proof and by property-based testing on reachable states. Given it,
  Theorem~\ref{thm:conv} is the ordinary canonicalization instantiated with
  Lemma~\ref{lem:swap}.
\item \textbf{Design/target.} Theorem~\ref{thm:conv} and the end-to-end hosting
  of the RGA (and Peritext) as conditioned instances are the framework's
  intended guarantee; their assembly rests on the in-progress step above. We have
  verified, on reachable states, that the RGA converges even structurally, so
  $\eqobs$ may there be taken as equality --- a simplification, not a soundness
  requirement.
\end{itemize}

\paragraph{Relation to prior methodologies.} Whole-state VC catalogues (Neem)
are the $\Inv=\app=\top$ collapse of this framework. The use of observational
equivalence rather than structural equality follows earlier work on
\emph{convergence modulo observable behaviour} (Peepul); the contribution here is
to combine it with applicability conditioning and to isolate the commutation
discipline (faithfulness + freshness) that makes state-dependent, re-anchoring
operations tractable, so that the author need only \emph{describe the sensible
operations}.

\end{document}
